\chapter{Introduction}
\label{ch:introduction}

\section{Motivation}
\label{sec:motivation}

Large language models (LLMs) are trained on internet-scale corpora that inevitably contain information
we would later like a model \emph{not} to retain, such as an individual's private details, copyrighted
passages, or knowledge that could facilitate harm. Once such data has shaped the model's weights, removing its
influence is far from trivial. Retraining the model from scratch on a cleaned corpus is the only method
that offers a clean guarantee, but for models with billions of parameters this is prohibitively expensive
in compute, time, and energy, and it must be repeated every time a new deletion request arrives. This
mismatch between the cost of retraining and the frequency of deletion requests has motivated the field of
\emph{machine unlearning}, a family of post-hoc procedures that aim to remove the influence of a specified
subset of the training data from an already-trained model, ideally recovering the behaviour of a model
that had never seen that data in the first place \citep{cao2015forget,bourtoule2021machine}.

The need is not merely academic. Data-protection regulation such as the European Union's General Data
Protection Regulation enshrines a ``right to erasure'', obliging data controllers to delete personal data
on request \citep{gdpr2016}; copyright holders increasingly demand the removal of protected text; and the
AI-safety community treats the removal of hazardous capabilities (for example, knowledge relevant to
bioweapons or cyberattacks) as a concrete mitigation lever \citep{li2024wmdp}. In all three settings,
unlearning is attractive precisely because it promises the effect of retraining at a fraction of the cost.
A recent survey positions unlearning as a pivotal component in the responsible life-cycle management of
LLMs \citep{liu2025rethinking}.

This promise, however, rests on a claim that is surprisingly difficult to verify, namely that the unlearned
model has genuinely \emph{forgotten} the target knowledge rather than merely learned to \emph{hide} it. The
central argument of this thesis is that the standard way of checking this claim, by inspecting the model's
outputs, is not enough, and that reading the model's \emph{internal} computation with tools from
mechanistic interpretability yields both a sharper diagnosis and a more trustworthy metric.

\section{How Unlearning Is Done, and the Method Families}
\label{sec:method-families}

Practical LLM unlearning is usually posed as a fine-tuning problem over two sets, a \emph{forget set}
\(D_f\) (the data whose influence is to be removed) and a \emph{retain set} \(D_r\) (data whose behaviour
should be preserved). The methods studied in this thesis are representative of the four dominant families,
which differ in \emph{what} they optimise and \emph{where} in the model they act.

\paragraph{Gradient-based methods.}
The most direct approach is \emph{gradient ascent} on the forget set, which maximises rather than minimises
the language-modelling loss on \(D_f\), pushing the model away from its memorised continuations
\citep{jang2022knowledge,yao2023llmunlearning}. Because unconstrained ascent quickly destroys general
capability, practical variants such as \emph{Gradient Difference} (GradDiff) combine ascent on \(D_f\) with
ordinary descent on \(D_r\), trading forgetting against utility. Gradient-based methods are simple and
model-agnostic but are known to be unstable and to risk catastrophic collapse.

\paragraph{Preference-based methods.}
\emph{Negative Preference Optimization} (NPO) reframes forgetting through the lens of preference
optimisation \citep{zhang2024npo}. It can be read as Direct Preference Optimization
\citep{rafailov2023dpo} with only negative examples. The forget responses are treated purely as dispreferred
completions, and the NPO loss is an adaptively weighted form of gradient ascent whose weight vanishes once a
sample is unlearned. This yields a bounded objective that provably approaches collapse exponentially more
slowly than plain gradient ascent, and is currently among the most effective forget/utility trade-offs.

\paragraph{Refusal / redirection methods.}
Rather than suppress the answer, \emph{redirection} methods teach the model to emit an innocuous
substitute. The ``I don't know'' (IDK) baseline fine-tunes the model to respond to forget queries with an
explicit expression of ignorance while preserving behaviour elsewhere \citep{maini2024tofu}. Such methods
change the model's \emph{output policy} on the forget distribution without necessarily disturbing the
underlying representation of the fact.

\paragraph{Representation-based methods.}
Finally, \emph{Representation Misdirection for Unlearning} (RMU) acts directly on the model's internal
activations \citep{li2024wmdp}. It steers the hidden states elicited by forget inputs towards random
vectors at a chosen set of layers, while regularising the retain-set activations to stay close to the
original model. RMU therefore intervenes on the computation itself rather than only on the output, which
makes it a natural stress-test for internal readouts.

These families make different implicit promises. Gradient ascent and NPO aim to lower the \emph{probability}
of the forget answer; IDK changes \emph{what is said}; RMU scrambles \emph{how the fact is represented}. A
key question of this thesis is whether an evaluation that only observes the output can tell these mechanisms
apart, and whether it should trust any of them as evidence of true erasure.

\section{The Verification Problem of Suppression versus Erasure}
\label{sec:verification-problem}

Unlearning is almost universally evaluated at the \emph{output}: one queries the model and checks whether
the target answer is still produced, using metrics such as question-answering accuracy, exact-match, ROUGE
overlap, the probability assigned to the gold answer, or membership-inference signals
\citep{maini2024tofu,shi2024muse,li2024wmdp}. These metrics are convenient and model-agnostic, but they
share a blind spot, in that they can only observe \emph{behaviour}, and behaviour can be changed without
removing the underlying knowledge.

A growing body of evidence shows that this blind spot is routinely exploited, whether deliberately or not.
Quantising an unlearned model back to lower precision restores much of the ``forgotten'' knowledge, with
retention rising from roughly a fifth to over four-fifths of its original level; the authors observe that
current benchmarks ``fail to detect'' this because they only measure outputs \citep{zhang2024quantization}.
A few steps of fine-tuning, in-context examples, jailbreak prompts, or a change of language can likewise
revive the target \citep{lynch2024eight,fan2025relearn}. Most tellingly for this thesis, when the
information is probed \emph{inside} the network rather than at the output, it is often plainly present:
logit-lens attacks recover ``deleted'' facts from the intermediate hidden states of edited models a large
fraction of the time \citep{patil2023candelete}. Independent work on semantic unlearning makes the same
diagnosis, noting that most methods operate at the lexical level and merely suppress a subset of surface
tokens, so that the forgotten knowledge ``remains accessible via paraphrasing'' \citep{semantic2025unlearning}.

The common thread is a gap between \emph{suppression} and \emph{erasure}. A metric that only sees the output
cannot distinguish a model that has genuinely removed a fact from one that has confidently learned to
deflect it, nor from one whose representation of the fact is intact but momentarily unexpressed. Closing
this gap requires a way to look inside.

\section{Reading Inside with Mechanistic Interpretability}
\label{sec:mech-interp}

Mechanistic interpretability studies the internal computation of a network rather than its input-output
behaviour alone. A useful organising picture is the \emph{residual stream}: a transformer maintains a
running representation to which each attention and feed-forward block reads and writes an additive update,
progressively refining a distribution over the next token \citep{elhage2021framework}. Factual knowledge in
particular has been localised to specific components, with feed-forward layers behaving like key-value
memories \citep{geva2021kv} and individual facts editable at identifiable mid-network sites
\citep{meng2022rome}.

To read what an intermediate state ``means'', one needs a \emph{lens}: a map from a hidden state at some
layer into the vocabulary, so that one can record the rank or probability the model would assign to the
correct answer at that depth. This thesis studies three lenses of increasing fidelity to the model's own
computation. The \emph{logit lens} applies the model's own final unembedding directly to an intermediate
state \citep{nostalgebraist2020logitlens}; it is parameter-free and fully reproducible, but noisy at early
layers. The \emph{tuned lens} learns a small affine correction per layer so that the readout reproduces the
final distribution \citep{belrose2023tunedlens}; it is better calibrated mid-stack but, being trained on the
output, tends to surface the answer earlier than the model actually computes it. The \emph{Jacobian lens}
transports a hidden state through the model's own averaged input-output Jacobian before decoding it
\citep{anthropic2026workspace}; it is \emph{interventionally grounded}, a first-order reflection of the
model's local input-output sensitivity rather than a fitted correlation, at a higher computational cost. A recurring theme of this
thesis is that these three readouts answer different questions, and that the choice of lens matters when one
wants to build a \emph{metric} rather than merely a visualisation.

\section{Thesis Question and Contributions}
\label{sec:contributions}

The benchmark that makes this study possible is Hubble \citep{allegro2025hubble}, an open suite of LLMs
(at 1B and 8B parameters) pretrained with a controlled set of facts inserted, so that for every fact one has
two ground-truth reference models. A \emph{Standard} model that never saw the fact (a lower bound on
knowledge) and a \emph{Perturbed} model trained to memorise it (an upper bound). Notably, Hubble untied its
output embeddings \emph{specifically} to support logit- and tuned-lens interpretability, and its authors
explicitly call for follow-up work that probes whether knowledge is localised and that develops more robust,
validated memorisation metrics; yet Hubble itself ships only output-level evaluation. This thesis takes up
that invitation.

The guiding question is therefore: \emph{can a mechanistic-interpretability readout of an unlearned model's
internal computation yield an evaluation metric that measures knowledge rather than surface form, and can we
validate that it does?} We decompose it into four research questions:
\begin{description}
  \item[\textbf{RQ1.}] Do output and lexical metrics distinguish the four unlearning method families, or do
  they collapse them onto a single ``forgotten'' verdict?
  \item[\textbf{RQ2.}] Does a layer-wise internal readout reveal residual internal accessibility that output
  metrics miss, and does it separate the four methods where the output cannot?
  \item[\textbf{RQ3.}] Can a metric derived from that readout be shown \emph{faithful}, that is, both
  knowledge-discriminative and paraphrase-stable, under the field's own meta-evaluation, and does
  it outperform established output and lexical metrics?
  \item[\textbf{RQ4.}] How do the three lenses (logit, tuned, Jacobian) differ as the basis for such a
  metric, and how does the picture change across model scale?
\end{description}
The contributions, which answer these questions in turn, are:

\begin{itemize}
  \item \textbf{A mechanistic addition to Hubble.} We integrate a corrected, boundary-safe logit lens into
  the Hubble evaluation, producing rank-by-layer trajectories for all six models (Standard, Perturbed, and
  four unlearned: NPO, GradDiff, IDK, RMU) at both scales, each anchored between the never-knew and knows
  references (Chapters~\ref{ch:introduction} onward).
  \item \textbf{A faithful metric, \(D_\text{mech}\).} From the logit-lens trajectory we derive a
  dual-anchor distance, \(D_\text{mech}\), that asks whether a fact's internal trajectory looks more like
  the ``knows'' or the ``never-knew'' reference, independent of what the model outputs.
  \item \textbf{Evidence that \(D_\text{mech}\) is faithful.} Under the field's own faithfulness
  meta-evaluation based on OpenUnlearning \citep{dorna2025openunlearning}, on Hubble-8B \(D_\text{mech}\) is
  the only evaluated metric that combines strong knowledge discrimination with substantial paraphrase
  stability (invariance shift \(\approx 0.47\) versus \(\geq 0.95\) for lexical and output metrics); its
  composite faithfulness score is about eighteen times the best lexical metric, an ordering robust to the
  choice of aggregation, and it replicates across a second benchmark (OpenUnlearning TOFU) at three scales.
  \item \textbf{A characterisation of all three lenses.} We compare the logit, tuned, and Jacobian lenses on
  the same facts, showing empirically why the tuned lens ``skips ahead'' and therefore cannot ground the
  metric, and implementing the first-order, interventionally grounded Jacobian lens end to end as a
  confirmatory readout.
  \item \textbf{A cross-scale method verdict.} The internal readout reveals a taxonomy invisible at the
  output. At 8B the lens cleanly separates a model that knows (rank~1) from one that never knew (rank~8),
  while the four unlearning methods leave distinct internal signatures, and this behaviour is strongly
  scale-dependent between 1B and 8B.
\end{itemize}

\section{Outline}
\label{sec:outline}

Chapter~\ref{ch:background} fixes notation and provides the technical background. The transformer and its
residual stream, the machine-unlearning problem, the loss functions of the four methods, the three lenses
in formal terms, and the definitions of the competing metrics. Chapter~\ref{ch:related-work} reviews the
literature on machine unlearning, its benchmarks and evaluation, the evidence that current methods suppress
rather than erase, and mechanistic interpretability, and situates the research gap.
Chapter~\ref{ch:methodology} describes our experimental design and the design of the \(D_\text{mech}\)
metric, including the faithfulness meta-evaluation protocol. Chapter~\ref{ch:results} reports the results,
answering RQ1 to RQ4 in turn. Chapter~\ref{ch:discussion} interprets them, relates them to prior work, and
states the threats to validity, and Chapter~\ref{ch:conclusion} concludes and outlines future work.
Additional detail, including hyperparameters, the Jacobian-lens derivation, supplementary results, and
reproducibility information, is provided in the appendices.
